\section[Aplicación del teorema de existencia de haces]
{Aplicación del teorema de existencia de haces: teorema de semicontinuidad
para los grupos fundamentales de las fibras de un morfismo propio y separable}
\label{X.2}
\marginpar{268}

\begin{theoreme}
\label{X.2.1}
\ifthenelse{\boolean{orig}}
{Sean~$Y$ el espectro de un anillo local noetheriano \emph{completo}, de
cuerpo residual~$\kres$,~$X$ un~$Y$-esquema propio,
$X_0=X\otimes_A\kres$,~$a_0$ un punto geométrico de~$X_0$ y~$a$ el punto
geométrico correspondiente de~$X$.}
{Sean~$A$ un anillo local noetheriano \emph{completo} de cuerpo residual~$k$,
$Y=\Spec(A)$,~$X$ un~$Y$-esquema propio, $X_0=X\otimes_Ak$,~$a_0$ un punto
geométrico de~$X_0$ y~$a$ el punto geométrico correspondiente de~$X$.}
Entonces el homomorfismo canónico~$\pi_1(X_0,a_0)\to\pi_1(X,a)$ es un
\emph{isomorfismo}.
\end{theoreme}

No es sino una traducción, al lenguaje del grupo fundamental, del resultado
recordado en~\Ref{IX.1.10}. Es aquí donde el teorema de existencia
\ifthenelse{\boolean{orig}}
{de}
{de}
haces en geometría algebraico-formal entra de manera esencial en la teoría del
grupo fundamental.

Introduzcamos ahora una clausura algebraica~$\bar{k}$ del cuerpo residual~$k$ y
la fibra geométrica~$\overline X_0=X_0\otimes_k\bar{k}$.
\ifthenelse{\boolean{orig}}
{}
{Elijamos un punto geométrico~$\overline a_0$ de~$\overline X_0$ sobre~$a_0$.}
Se tiene, pues, la sucesión exacta (\Ref{IX.6.1})
$$
\ifthenelse{\boolean{orig}}
{e\to\pi_1(\overline X_0,\overline a_0)\to\pi_1(X_0,a_0)\to
\pi_1(k,\bar{k})\to e,}
{e\to\pi_1(\overline X_0,\overline a_0)\to\pi_1(X_0,a_0)\to
\pi_1(k,\bar{k})\to e,}
$$
por otra parte, se tienen el isomorfismo~\Ref{X.2.1} y el isomorfismo análogo,
más elemental,
$$
\pi_1(k,\bar{k})\to\pi_1(Y,b),
$$
donde~$b$ es la imagen de~$a$ en~$Y$. Se obtiene así:

\begin{corollaire}
\label{X.2.2}
\ifthenelse{\boolean{orig}}
{Con las notaciones anteriores, supongamos conexo~$\overline X_0$ y sean
$\overline a_0$ un punto geométrico de
$\overline X_0=X_0\otimes_k\bar{k}$,~$a_0$ su imagen en~$X$,}
{Con las notaciones anteriores, supongamos conexo
$\overline X_0=X_0\otimes_k\bar{k}$ y sean~$\overline a_0$ un punto geométrico
de~$\overline X_0$,~$a_0$ su imagen en~$X$,}
~$b_0$ su imagen en~$Y$. Entonces la sucesión siguiente de homomorfismos
canónicos
$$
\ifthenelse{\boolean{orig}}
{e\to\pi_1(\overline X_0,\overline a_0)\to\pi_1(X,a_0)\to
\pi_1(Y,b_0)\to e}
{e\to\pi_1(\overline X_0,\overline a_0)\to\pi_1(X,a_0)\to
\pi_1(Y,b_0)\to e}
$$
es exacta.
\end{corollaire}

Compárese esta sucesión con la sucesión exacta~\Ref{X.1.4}, pero obsérvese que
a)~no ha sido necesario imponer ninguna hipótesis de planitud ni de
separabilidad sobre las fibras para~$X\to Y$; y b)~se obtiene el complemento
importante de que \emph{el morfismo
$\pi_1(\overline X_0,\overline a_0)\to\pi_1(X,a_0)$ es inyectivo}.

Este último hecho permitirá comparar el grupo fundamental de las demás
\marginpar{269}
fibras geométricas de~$X$ sobre~$Y$ con el de~$\overline X_0$. Sea, en efecto,
~$y_1$ un punto cualquiera de~$Y$,~$X_1$ su fibra y~$\overline X_1$ su fibra
geométrica respecto de una extensión algebraicamente cerrada de~$\kres(y_1)$.
Sea~$\overline a_1$ un punto geométrico de~$\overline X_1$,~$a_1$ su imagen
en~$X$ y~$b_1$ su imagen en~$Y$. Elijamos una «clase de caminos» de~$a_1$
a~$a_0$, de donde se obtiene una clase de caminos de~$b_1$ a~$b_0$ y, por
tanto, un diagrama conmutativo de homomorfismos:
$$
\xymatrix@=.5cm{
&\pi_1(\overline X_1,\overline a_1) \ar[r] & \pi_1(X,a_1) \ar[r]\ar[d] & \pi_1(Y,b_1) \ar[r]\ar[d] & e \\
e \ar[r]& \pi_1(\overline X_0,\overline a_0) \ar[r] & \pi_1(X,a_0) \ar[r] &\pi_1(Y,b_0) \ar[r] & e,
}
$$
en el cual las dos flechas verticales indicadas son isomorfismos. Como la
segunda fila es exacta, se obtiene un homomorfismo canónico, que llamaremos
\emph{homomorfismo de especialización para el grupo fundamental} (que solo
depende de la clase de caminos elegida de~$a_1$ a~$a_0$ y está, por tanto,
\emph{definido módulo un automorfismo interior de}~$\pi_1(X,a_0)$):
$$
\pi_1(\overline X_1,\overline a_1)\to
\pi_1(\overline X_0,\overline a_0).
$$
Cuando la primera fila anterior es también exacta, se sigue inmediatamente que
el homomorfismo de especialización es sobreyectivo. Teniendo en cuenta
~\Ref{X.1.4}, se obtiene, pues:

\begin{corollaire}
\label{X.2.3}
Bajo las hipótesis de~\Ref{X.2.1}, supongamos además que el morfismo
~$f\colon X\to Y$ es \emph{separable} \eqref{X.1.1} y que~$\overline X_0$ es
conexo (por tanto, en virtud de~\Ref{X.1.2}, se tiene
$f_*(\cal{O}_X)=\cal{O}_Y$). Entonces, para toda fibra geométrica
~$\overline X_1$ de~$X$ sobre~$Y$, provista de un punto geométrico
~$\overline a_1$, el homomorfismo de especialización
\index{homomorfismo de especialización para el grupo fundamental|hyperpage}%
\index{especialización para el grupo fundamental (homomorfismo de)|hyperpage}%
definido anteriormente es \emph{sobreyectivo}.
\end{corollaire}

Este es un resultado de \emph{semicontinuidad} para el grupo fundamental, que
no parece tener todavía un análogo en topología algebraica. Puede enunciarse,
por lo demás, bajo hipótesis más generales:

\begin{corollaire}
\label{X.2.4}
Sean
\marginpar{270}
~$f\colon X\to Y$ un morfismo propio de fibras geométricamente conexas,
con~$Y$ localmente noetheriano;~$y_0$ e~$y_1$ dos puntos de~$Y$ tales que
$y_0\in\overline{\{y_1\}}$;~$\overline X_0$ y~$\overline X_1$ las fibras
geométricas de~$X$ correspondientes a extensiones algebraicamente cerradas
elegidas de~$\kres(y_0)$ y~$\kres(y_1)$; y~$\overline a_0$,
\resp~$\overline a_1$, un punto geométrico de~$\overline X_0$,
\resp~$\overline X_1$. Entonces puede definirse de manera natural un
homomorfismo de especialización
$$
\pi_1(\overline X_1,\overline a_1)\to
\pi_1(\overline X_0,\overline a_0),
$$
definido salvo automorfismo interior, y este homomorfismo es \emph{sobreyectivo}
si~$f$ es un morfismo separable \eqref{X.1.1}.
\end{corollaire}

En efecto, de~\Ref{X.1.8} se sigue primero que~\Ref{X.2.4} es esencialmente
independiente de las extensiones algebraicamente cerradas elegidas para los
cuerpos residuales~$\kres(y_0)$ y~$\kres(y_1)$. Esto permite sustituir~$Y$ por
un preesquema~$Y'$ sobre~$Y$ que tenga un punto~$y_0'$ (\resp~$y_1'$) sobre
~$y_0$ (\resp~$y_1$). Se toma entonces por~$Y'$ el espectro del completado del
anillo local de~$y_0$ en~$Y$, y se aplica~\Ref{X.2.3}.

\begin{remarques}
\label{X.2.5}
La conclusión final de~\Ref{X.2.4} sobre la sobreyectividad del homomorfismo de
especialización, y a fortiori los resultados~\Ref{X.1.3} y~\Ref{X.1.4} de los
que se deduce, deja de ser cierta si ya no se supone que~$f\colon X\to Y$ es
un morfismo separable, incluso para esquemas proyectivos sobre un cuerpo
algebraicamente cerrado de característica~$0$. Más adelante veremos ejemplos
tanto cuando~$f$ es plano pero~$f$ admite una fibra no separable (aunque~$X$
e~$Y$ sean lisos sobre~$k$), como cuando las fibras de~$f$ son separables pero~$f$ no
es plano (por ejemplo, cuando~$f\colon X\to Y$ es un morfismo birracional de
esquemas íntegros normales); véase~\Ref{XI.3}. En estos ejemplos puede
ocurrir que el grupo fundamental de la fibra geométrica genérica sea trivial,
pero no el de una fibra geométrica especial conveniente.

Por otra parte, incluso si~$f\colon X\to Y$ es un morfismo propio y separable
como en~\Ref{X.2.4}, es frecuente que el homomorfismo de especialización no sea
\marginpar{271}
un isomorfismo. Por ejemplo, es fácil dar casos en los que~$\overline X_1$ es
una curva elíptica no singular (por tanto, su grupo fundamental es conmutativo
y su componente~$\ell$-primaria, para un número primo~$\ell$ distinto de la
característica, es isomorfa a~$\ZZ_\ell^2$; véase~\Ref{XI}), mientras que
~$\overline X_0$ está formada ya por dos curvas racionales no singulares que se
cortan en dos puntos, ya por dos curvas racionales tangentes en un punto, ya,
por último, por una curva racional con un punto singular que es una cúspide.
(Para la clasificación completa de las curvas elípticas degeneradas, véanse
los trabajos recientes de Kodaira~\cite{X.1} y Néron). En estos casos, el grupo
fundamental de~$\overline X_0$ es, respectivamente,~$\hat\ZZ$,~$e$,~$e$, y es,
por tanto, «estrictamente menor» que el de~$\overline X_1$.

Veremos, sin embargo, más adelante, cuando~$f$ es un morfismo \emph{liso}, una
cota para el núcleo del homomorfismo de especialización que implica, en
particular, que si~$\kres(y_0)$ tiene característica~0, el homomorfismo de
especialización es un isomorfismo. Pero incluso para un morfismo liso, si la
característica de~$\kres(y_0)$ es~$>0$, el homomorfismo de especialización
puede no ser un isomorfismo, como se ve, por ejemplo, cuando~$X$ es un esquema
abeliano sobre~$Y$ (incluso de dimensión relativa~$1$),
\ifthenelse{\boolean{orig}}
{\cf~\Ref{XI.2}).}
{\cf~\Ref{XI.2}.}
Una teoría satisfactoria de la especialización del grupo fundamental debe
tener en cuenta la «componente continua» del «verdadero» grupo fundamental,
correspondiente a la clasificación de los recubrimientos principales cuyos
grupos estructurales son grupos infinitesimales. Con ello cabría esperar que
los «verdaderos» grupos fundamentales de las fibras geométricas de un morfismo
liso y propio~$f\colon X\to Y$ formasen un sistema local bien comportado
sobre~$X$, límite proyectivo de esquemas en grupos finitos y planos sobre~$X$
\kern1pt\footnote{Esta conjetura sumamente seductora queda, por desgracia,
refutada por un ejemplo inédito de M\ptbl Artin, ya cuando las fibras de~$f$
son curvas algebraicas de un género dado~$g\ge 2$.}. Volveremos más adelante
sobre este punto de vista; nuestro propósito presente es, por el contrario,
llevar tan lejos como sea posible los fenómenos comunes a la teoría topológica
y a la teoría esquemática del grupo fundamental.
\end{remarques}

Sea ahora~$X_0$ una curva propia, lisa y conexa de género~$g$ sobre un cuerpo
algebraicamente cerrado~$k$. Si~$k$ tiene característica cero, su grupo
fundamental puede determinarse por métodos trascendentes como sigue. Se sabe
que~$X_0$ procede por extensión de la base de una curva definida
\marginpar{272}
sobre una extensión algebraicamente cerrada de grado de trascendencia finito
del cuerpo primo~$\QQ$ y, teniendo en cuenta~\Ref{X.1.8}, puede suponerse que
~$k$ mismo tiene grado de trascendencia finito sobre~$\QQ$. Puede suponerse,
por tanto, que~$k$ es un subcuerpo del cuerpo~$\CC$ de los números complejos,
y una nueva aplicación de~\Ref{X.1.8} permite suponer que~$k=\CC$. No es
difícil comprobar entonces que el grupo fundamental de
\ifthenelse{\boolean{orig}}
{$X$}
{$X_0$}
es isomorfo a la completación profinita del grupo fundamental del espacio
topológico asociado
\ifthenelse{\boolean{orig}}
{$\tilde X$}
{$\widetilde X_0$}
(una superficie compacta orientada de género~$g$), para la topología definida
por los subgrupos de índice finito\footnote{Esta deducción se explicó en una de
las exposiciones orales que no fueron redactadas.}. Por otra parte, es clásico
que el grupo fundamental topológico está engendrado por~$2g$ generadores
~$s_i,t_i$ ($1\le i\le g$), sujetos a la única relación
$$
(s_1t_1s_1^{-1}t_1^{-1})\cdots(s_gt_gs_g^{-1}t_g^{-1})=1.
$$
Por tanto, el grupo fundamental de
\ifthenelse{\boolean{orig}}
{$X$}
{$X_0$}
admite~$2g$ generadores \emph{topológicos}~$s_i,t_i$ ($1\le i\le g$), ligados
por la única relación anterior.

Supongamos ahora que la característica de~$k$ es~$p>0$. Sea~$A$ el anillo de
vectores de Witt construido a partir de~$k$, y~$K$ una extensión
algebraicamente cerrada de su cuerpo de fracciones. Se vio en
~\Ref{III.7.4} que existe un esquema~$X$ sobre~$Y=\Spec(A)$, propio y liso
sobre~$Y$, cuya reducción es~$X_0$. Aplicando~\Ref{X.2.3}, se obtiene un
homomorfismo \emph{sobreyectivo}
$$
\pi_1(X_1)\to\pi_1(X_0),
$$
donde~$X_1=X\otimes_AK$. Es inmediato\footnote{\Cf EGA~IV 12.2.} que~$X_1$
es liso sobre~$K$, conexo por~\eqref{X.1.2}, de dimensión~$1$, y que su género
es igual a~$g$ por la invariancia de la característica de Euler--Poincaré;
véase EGA~III 7. Como~$K$ tiene característica~$0$, puede aplicársele el
resultado anterior. Se ha probado así, por \emph{métodos trascendentes}:

\begin{theoreme}
\label{X.2.6}
Sea~$X_0$ una curva algebraica lisa, propia y conexa sobre un
\marginpar{273}
cuerpo algebraicamente cerrado~$k$, y sea~$g$ su género. Entonces
~$\pi_1(X_0)$ admite un sistema de~$2g$ generadores topológicos, ligados por la
relación escrita anteriormente. Si~$k$ tiene característica~$0$,
~$\pi_1(X_0)$ es además el grupo libre de tipo galoisiano sobre esos generadores
sujeto a dicha relación.
\end{theoreme}

\begin{remarques}
\label{X.2.7}
En la actualidad, que sepa el redactor, no existe una demostración por métodos
puramente algebraicos del resultado anterior, salvo para los géneros~$0,1$.
Para empezar, no se ve en absoluto cómo distinguir en~$\pi_1(X_0)$
~$2g$ elementos de los que después pudiera esperarse que formasen un sistema de
generadores topológicos. A este respecto, es más favorable la situación de la
recta racional privada de~$n$ puntos y el estudio de sus recubrimientos
moderadamente ramificados en esos puntos, pues los grupos de ramificación en
los~$n$ puntos proporcionan~$n$ elementos del grupo fundamental estudiado y se
prueba que lo engendran topológicamente, como veremos más adelante
\footnote{\Cf Exposición~\Ref{XII}. Incluso estos elementos solo quedan
determinados módulo conjugación, y conviene elegir \emph{simultáneamente} y con
cuidado representantes de sus clases.}. Pero ni siquiera en este caso
particularmente concreto parece conocerse una demostración puramente
algebraica. Una demostración tal sería, desde luego, sumamente interesante. El
único hecho relativo al grupo fundamental de una curva que parece admitir una
demostración puramente algebraica, aparte del teorema débil de finitud
~\Ref{X.2.12} que Lang--Serre~\cite{X.2} probaron algebraicamente, es la
determinación de su abelianización mediante la jacobiana, señalada en la última
línea de~\Ref{IX.5.8}.
\end{remarques}
